\section{Extensions: Aliasing, Mesh Heterogeneity, and Parameter Decomposition}
\label{sec:ext}

The validation above used the default linear-Gaussian generator and a CNN
forecaster, and several GraphCast-relevant questions cannot even be posed in that
setting. Discovery ran at the cadence the dynamics are defined on, so nothing ever
aliased. A translation-equivariant CNN has no per-position parameters and no
privileged locations, so grid-locked features cannot arise in it, whatever the
data. And linear-Gaussian dynamics leave contemporaneous directions unidentifiable
in principle while pinning the forecaster's representation to a nearly low-rank
subspace. The experiments in this section relax each of these constraints.

\subsection{A harder generator}
\label{sec:finecadence}

The fine-cadence generator (\texttt{generate\_finecadence.py}; 100 realisations,
$T=2400$) keeps the eight-mode edge set of the default graph and changes the rest.
Cross-edges get heterogeneous lags $\ell \in \{1,2,3,4,6\}$ in fine-step units. The
dynamics gain a saturating autoregression and a bilinear interaction term
($\alpha=0.5$, $\beta=0.15$). Innovations become skew-normal (shape $\alpha=4$,
standardized to unit variance, so VAR stationarity is untouched). The skewness is
there for a reason: it is the LiNGAM condition under which contemporaneous
directions become identifiable in principle (Sec.~\ref{sec:background}), so this
generator can tell us whether our discovery methods actually cash in on that
identifiability.

\subsection{Subsampling and aliased edges: PCMCI+ vs.\ PCMCI}
\label{sec:subsample}

Observing the fine-cadence system at stride $s$ turns a ground-truth lag $\ell$
into a coarse lag $\lfloor \ell/s \rfloor$, so every edge with $\ell < s$ aliases
to a contemporaneous ($\tau=0$) dependency. Plain PCMCI with $\tau_{\min}=1$ has no
way to represent such an edge; its recall on aliased edges is zero before it sees
any data. Table~\ref{tab:subsample} shows the sweep over strides for PCMCI+
($\tau_{\min}=0$) and PCMCI ($\tau_{\min}=1$) on the latent modes (40 realisations,
$\alpha=0.05$).

\begin{table}[h]
\centering
\small
\begin{tabular}{cccccc}
\toprule
$s$ & $T$ & GT edges (lag/contemp.) & PCMCI+ contemp.\ F1 & PCMCI+ overall F1 &
PCMCI overall F1 \\
\midrule
1 & 2400 & 12 / 0 & ---   & 0.776 & 0.593 \\
2 & 1200 &  9 / 3 & 0.627 & 0.587 & 0.457 \\
3 &  800 &  6 / 6 & 0.491 & 0.444 & 0.308 \\
4 &  600 &  4 / 8 & 0.376 & 0.346 & 0.192 \\
\bottomrule
\end{tabular}
\caption{Subsampling sweep on the fine-cadence dataset. ``GT edges'' counts
ground-truth edges that remain lagged vs.\ alias to $\tau=0$ at stride $s$.
Contemporaneous F1 is computed on the undirected adjacency. At $s=1$ the PCMCI+
lagged F1 is $0.825$, matching the default-dataset ceiling of
Sec.~\ref{sec:savar}.}
\label{tab:subsample}
\end{table}

PCMCI+ recovers a substantial fraction of the aliased edges (contemporaneous F1
$0.38$--$0.63$) while plain PCMCI's overall F1 collapses with aliasing, $0.593$
down to $0.192$. Whenever dynamics might act faster than the sampling
interval---and at 6-hour steps on the real atmosphere they certainly do---PCMCI+
is the tool. Its remaining weakness is orientation. Most recovered $\tau=0$ edges
come back unoriented (orientation rate $0.08$--$0.30$), even though the skewed
innovations make the directions identifiable in principle. ParCorr is a Gaussian
test statistic; the directional information lives in higher moments it never
looks at. Bolting a LiNGAM-style orientation step onto PCMCI+'s adjacencies is
the obvious follow-up and has not been done yet (Sec.~\ref{sec:open}).

\subsection{A GraphCast-faithful forecaster}
\label{sec:gnn}

For grid-locked features to be possible at all, the forecaster needs privileged
locations. We replaced the CNN with a message-passing GNN
(\texttt{train/gnn\_forecaster.py}) built like GraphCast's processor in miniature:
a mesh over the $50\times50$ grid with local 8-neighbour edges everywhere plus
long-range edges among a lattice of hub nodes, node degrees running from 3 to 32,
and message-passing MLPs shared across all nodes. The mesh is the only source of
positional asymmetry. The SAVAR generator supplies none: the mode layout is
content, tied to the causal graph, and the observation noise is spatially
homogeneous. Any hub-tied structure in the trained model is therefore attributable
to the mesh.

An early version of this model also had a learnable per-node embedding, added on
the reasoning that a CNN lacks per-position parameters and the GNN should have
them. We removed it. GraphCast's processor has no such lookup table, and giving
position a parametric home of its own would have manufactured the grid-lock
phenomenon we were trying to observe emerging from the mesh. All results below use
the model without it.

The trained GNN reaches validation RMSE $0.4295$ on the fine-cadence one-step
task, against $0.4591$ for predicting the grand mean: $R^2 = 0.125$. This is a
weak-signal task by construction (the one-step ceiling at the mode level is
$R^2 \approx 0.20$; Sec.~\ref{sec:homog}), so throughout this section we read the
structure of the interpretability results and put little weight on absolute effect
sizes.

\subsection{Per-mode SAEs on the GNN}
\label{sec:gnn-sae}

We repeated the per-mode SAE pipeline of Sec.~\ref{sec:savar} on the GNN's node
states after the final message-passing layer (mode-weighted pooling, eight
independent TopK SAEs, ridge ceilings measured out-of-fold). The picture is the
CNN's, only sharper. Alignment drops from $7/8$ modes to $2/8$ (X1 and X2), still
with no strong alignment and no monosemantic feature anywhere. The global-activity
direction that dominated the CNN's pooled activations at ${\approx}86\%$ of
variance dominates the GNN's at $88$--$92\%$, and the linear decoding ceilings
shrink accordingly ($0.36$--$0.50$, versus $0.36$--$0.59$ for the CNN). The SAEs
themselves are not the bottleneck; they still reach $66$--$81\%$ of ceiling. No
feature is mode-specific (specificity $\le +0.07$), and modes X1 and X6 literally
share their best feature. Message passing over the heterogeneous mesh concentrates
node state onto global system activity even more aggressively than convolution
does, and mode-specific structure never surfaces in the representation. The
parameter-level analysis below says something similar about the hub structure.

\subsection{Parameter decomposition of the message-passing MLPs}
\label{sec:vpd}

We applied VPD to the frozen GNN: both linear maps of all four message-passing
MLPs, $C=64$ rank-one components each, per-node gates, stochastic plus adversarial
(PGD) masks, frequency-minimality on. Reconstruction converged: masked-forward
MSE to the frozen model's output of ${\approx}2\times10^{-3}$ under stochastic
masks, ${\approx}3.5\times10^{-2}$ under worst-case masks, faithfulness residual
${\approx}8\times10^{-4}$. The decomposition reproduces the forecaster even under
worst-case node-level ablations, which is the property that licenses reading the
gates causally.

Each component's gate map (mean gate per node, reshaped to $50\times50$) is
summarized by its \emph{hub contrast} (mean gate on hub nodes minus background),
its \emph{mode contrast} (mean gate on the eight mode blobs minus background), and
an ablation effect (rise in reconstruction error when the component is fully
masked). Table~\ref{tab:vpd} gives per-layer means.

\begin{table}[h]
\centering
\small
\begin{tabular}{lccccc}
\toprule
Layer & Hub contrast & Mode contrast & Abl.\ mean ($\times10^{-3}$) &
Abl.\ max ($\times10^{-3}$) & Mode-dominant \\
\midrule
0, MLP in   & \textbf{0.479} & 0.369 & 1.86 & 17.3 & 19/64 \\
0, MLP out  & \textbf{0.258} & 0.113 & 0.62 &  2.1 & 18/64 \\
1, MLP in   & \textbf{0.441} & 0.327 & 3.08 & 37.1 & 21/64 \\
1, MLP out  & \textbf{0.222} & 0.149 & 0.20 &  0.6 & 26/64 \\
2, MLP in   & 0.151 & \textbf{0.340} & 2.40 & 25.1 & 50/64 \\
2, MLP out  & 0.006 & \textbf{0.227} & 0.10 &  1.2 & 48/64 \\
3, MLP in   & 0.118 & \textbf{0.365} & 1.24 &  8.2 & 48/64 \\
3, MLP out  & 0.041 & 0.049 & 0.03 &  0.9 & 41/64 \\
\bottomrule
\end{tabular}
\caption{VPD gate structure per decomposed matrix (64 components each).
``Mode-dominant'' counts components whose gate is at least as mode-locked as
hub-locked. Bold marks the larger of the two contrasts.}
\label{tab:vpd}
\end{table}

Depth organizes the gates. In layers 0--1 they concentrate on the mesh hubs: hub
contrast exceeds mode contrast, and only 19--26 of 64 components are more mode-
than hub-locked. By layers 2--3 the pattern has inverted, with 48--50 of 64
components tracking the SAVAR mode blobs. Early mechanisms lock to the mesh, late
mechanisms to content. Nearly all of this structure sits in the first sublayer of
each MLP, whose gate input includes the neighbour-aggregated state; that is enough
to tell a hub from an ordinary cell without any positional encoding.

The gates also settle the grid-lock question, this time from inside a single
object. Across early-layer components, hub contrast correlates negatively with
ablation effect: $-0.48$ at layer 0, $-0.37$ at layer 1. The most hub-locked
mechanisms are the ones the forecast depends on least. Section~\ref{sec:gridlock-cnn}
reached the same conclusion for the CNN's SAE features, but needed a separate
ablation experiment to get there; here the gate is ablation-defined, so
``position-locked'' and ``causally inert'' are read off the same map. No time
axis is involved either, so the aliasing issues of Sec.~\ref{sec:subsample} never
arise. For grid-lock specifically, the parameter route is simply the better
instrument.

Caveats: the ablation effects are small in absolute terms everywhere (means of
$1$--$3\times10^{-3}$), as expected against a weak-signal target, so the depth
gradient and the hub--ablation correlation are claims about structure, and we make
no claim about magnitudes. The components are also far from specialized: every
component's gate covers all eight mode blobs almost equally (coefficient of
variation ${\approx}0.05$). Why nothing specializes is the next question.

\subsection{Why nothing specializes}
\label{sec:homog}

The GNN's margin over the grand-mean predictor is thin ($R^2 = 0.125$), and its
behaviour is essentially damped persistence: correlation $0.78$ with the last
input frame, effective gain ${\approx}0.27$. The generator explains this. The
one-step mode-forecast ceiling is only $R^2 \approx 0.20$, and the eight modes are
dynamically near-interchangeable (variance ratio $1.1$, self-coupling
coefficients all within $[0.30, 0.55]$, cross-coupling worth a mere $0.03$ of
$R^2$). For such a system the optimal one-step forecaster is close to one shared
shrinkage toward the mean, applied identically everywhere. Mode-specific
circuitry would buy nothing, and the model built none; the SAE features of
Sec.~\ref{sec:gnn-sae} and the VPD components of Sec.~\ref{sec:vpd} are
homogeneous because the task is.

This account makes a testable prediction: give the modes genuinely different
dynamics and specialization should appear. A modified generator
(\texttt{generate\_hetdynamics.py}) widens the self-coupling range to
$[0.15, 0.92]$, a ${\sim}23\times$ spread in autocorrelation timescale, with an
equal-variance control to separate timescale heterogeneity from amplitude. The
three-way comparison (baseline, heterogeneous dynamics, equal-variance control)
through the full GNN and SAE pipeline is running at the time of writing.
\todo{Add heterogeneous-dynamics results when the runs complete.}
